\documentclass[11pt,a4paper]{article}

% ------------------------------------------------------------------ packages
\usepackage[utf8]{inputenc}
\usepackage[T1]{fontenc}
\usepackage{lmodern}
\usepackage{microtype}
\usepackage[margin=2cm]{geometry}
\usepackage{amsmath}
\usepackage{booktabs}
\usepackage{tabularx}
\usepackage{array}
\usepackage{enumitem}
\usepackage[numbers,sort&compress]{natbib}
\usepackage[hidelinks]{hyperref}

\setlist{nosep, leftmargin=*}
\setlength{\parskip}{0.4em}
\setlength{\parindent}{0pt}
\newcolumntype{L}{>{\raggedright\arraybackslash}X}
\newcommand{\var}[1]{\texttt{#1}}

% ------------------------------------------------------------------ title
\title{\textbf{Hidden Suicide Risk:} Can routinely recorded health information
       identify suicidal ideation that a depression screen misses?}
\author{Dimi \and Rana \and Roger Bavibidila}
\date{Data Science in Health Care --- Preliminary study}

\begin{document}
\maketitle

% =================================================================== 1
\section{Problem statement}

Suicide is a leading cause of preventable death. Switzerland records about 1,000
suicides per year (996 in 2024, 72\% men)~\cite{bfs2026}. Many people who die by
suicide were in contact with a physician in the preceding year~\cite{ahmedani2014}.
The routine tool for spotting mental distress in primary care is a depression
questionnaire such as the PHQ-9, with a positive screen at a score of 10 or
more~\cite{kroenke2001}. \textbf{But suicidal thoughts are not limited to people who
screen positive for depression.} In our pilot analysis of NHANES, \textbf{42\% of
adults who reported thoughts of self-harm or of being better off dead scored below 10
on the other eight PHQ items}, so a standard depression screen would miss them.

At the same time, an electronic health record already holds structured information a
clinician rarely has time to cross-reference: prescriptions (antidepressants, opioids,
benzodiazepines), healthcare use (hospital stays, mental-health visits), chronic
conditions, sleep problems, employment and insurance status. \textbf{It is unclear
whether this routinely recorded information carries enough signal to flag suicidal
ideation, especially in patients the depression screen does not catch.} This project
quantifies that signal on real, nationally representative health data.

% =================================================================== 2
\section{Research question and hypotheses}

\textbf{Main research question}
\begin{quote}
\itshape To what extent can routinely recorded health information identify adults with
current suicidal ideation, and does it add value beyond a depression screen, in
particular for adults who screen negative (PHQ-8~$<$~10)?
\end{quote}

\textbf{Sub-questions}
\begin{enumerate}
  \item How should the severe class imbalance (about 1:26) be handled: class
        weighting, SMOTE, or no correction?
  \item Does performance hold when a model trained on older survey waves is applied to
        newer ones, and is it similar across sex, age groups and ethnicity?
\end{enumerate}

\textbf{Hypotheses} (all tested on the 2015--2018 test set with paired bootstrap 95\% CIs)

\textit{Main question}
\begin{itemize}
  \item \textbf{H1:} A model using routine health data achieves a PR-AUC at least
        \textbf{twice the outcome prevalence}, and higher than a demographics-only model
        (CI of the difference excludes 0).
  \item \textbf{H2:} Adding routine health data to the PHQ-8 score increases PR-AUC
        compared with PHQ-8 alone (CI of the difference excludes 0).
  \item \textbf{H3:} Among adults who screen negative on depression (PHQ-8~$<$~10), the
        routine-data model ranks people with suicidal ideation above those without
        better than chance (ROC-AUC CI lies above 0.5), and its PR-AUC is at least twice
        the prevalence within this group.
\end{itemize}

\textit{Sub-question 1 (class imbalance)}
\begin{itemize}
  \item \textbf{H4:} Class weighting and SMOTE do not improve PR-AUC over the
        uncorrected model, but they worsen calibration (higher Brier score; CI of the
        difference excludes 0), because they shift predicted probabilities away from
        the true prevalence~\cite{vandengoorbergh2022}.
\end{itemize}

\textit{Sub-question 2 (robustness)}
\begin{itemize}
  \item \textbf{H5:} The model trained on 2005--2014 keeps its discrimination on
        2015--2018 (test ROC-AUC within 0.05 of the cross-validated ROC-AUC), and
        ROC-AUC differs by less than 0.05 between sexes, age groups and ethnic groups.
        ROC-AUC is used here rather than PR-AUC because PR-AUC depends on prevalence,
        which differs between waves and subgroups.
\end{itemize}

% =================================================================== 3
\section{Data sources}

\begin{table}[h]
\centering
\small
\begin{tabularx}{\textwidth}{@{}>{\raggedright\arraybackslash}p{3.1cm} L
                               >{\raggedright\arraybackslash}p{3.6cm}
                               >{\raggedright\arraybackslash}p{2.6cm}@{}}
\toprule
\textbf{Source} & \textbf{Content} & \textbf{Size / years} & \textbf{Role} \\
\midrule
\textbf{NHANES} (US CDC / NCHS)~\cite{nhanes}
  & Individual-level interview and examination data, linked by the respondent ID
    \var{SEQN}
  & 7 cycles 2005--2018 $\rightarrow$ \textbf{36,259 adults} with a complete PHQ-9,
    \textbf{1,366} with suicidal ideation (3.8\%)
  & Main modelling dataset \\
\addlinespace
NHANES modules used
  & \var{DPQ} (PHQ-9), \var{DEMO} (age, sex, ethnicity, education, income ratio,
    marital status, survey weights), \var{RXQ\_RX} (prescription drug names),
    \var{HUQ} (self-rated health, hospital stay, mental-health visit),
    \var{MCQ}/\var{DIQ}/\var{BPQ} (chronic conditions), \var{SLQ} (sleep trouble),
    \var{SMQ} (smoking), \var{DUQ} (drug use), \var{OCQ} (employment),
    \var{HIQ} (insurance), \var{BMX} (BMI)
  & 13 modules $\times$ 7 cycles $=$ 91 files
  & Feature construction \\
\addlinespace
NHANES 2021--2023 (\var{DPQ\_L})
  & PHQ-9 only (the public drug file lacks drug names)
  & 5,455 adults
  & Descriptive trend check only \\
\addlinespace
\textbf{Swiss Federal Statistical Office (BFS)} via opendata.swiss~\cite{bfs2026}
  & Suicides by method, sex and age (1995--2024); assisted suicide by sex and age
    (2003--2024)
  & Aggregate tables
  & Swiss context and motivation (not modelled) \\
\bottomrule
\end{tabularx}
\end{table}

\textbf{Outcome:} \var{si} $= 1$ if PHQ-9 item 9 (``thoughts that you would be better
off dead, or of hurting yourself'') is answered on at least \textit{several days} in
the past two weeks. A stricter outcome (\textit{more than half the days}, 424 cases,
$\approx$1.2\%) is used as a sensitivity analysis.

\textbf{Leakage rule:} PHQ items 1--8 are \textbf{excluded} from the ``routine data''
feature set and enter only as the PHQ-8 score in the explicit comparison (H2).

% =================================================================== 4
\section{Approach and methodology}

\textbf{Step 1 --- Acquisition and harmonisation (week 1, done).} Download the 91
\var{.xpt} files, verified against SHA-256 checksums, and merge all modules on
\var{SEQN} per cycle. Restrict to adults (18+) with a complete PHQ-9. SAS transport
stores 0 as $\approx 5\cdot10^{-79}$, so these values are set to 0 on import.
``Refused / don't know'' codes are recoded to missing per variable. Variables that
change between cycles are harmonised, for example \var{HUQ050} $\rightarrow$
\var{HUQ051} (healthcare visits, 2013) and the new sleep module from 2015; each
variable is kept only for the ages asked in every cycle. All rules are listed in
Appendix~\ref{app:harmonisation}. Outputs: one row per adult, prescriptions in long
format (one row per drug), the 2021--23 PHQ-9 file, and tables of sample size and
missingness per cycle.

\textbf{Step 2 --- Feature engineering (week 1--2).}
\begin{itemize}
  \item \textit{Medication:} map free-text drug names in \var{RXQ\_RX} (one row per
        drug) to four classes (antidepressants, opioids, benzodiazepines/Z-drugs,
        antipsychotics/mood stabilisers) plus a \textbf{polypharmacy} count of
        distinct drugs.
  \item \textit{Multimorbidity:} count of 10 diagnosed chronic conditions (arthritis,
        heart disease, stroke, COPD, cancer, diabetes, hypertension, \ldots).
  \item \textit{Healthcare use:} hospital stay and mental-health visit in the past
        year, number of visits, self-rated health.
  \item \textit{Social:} poverty-income ratio, not working, uninsured, marital status,
        education.
  \item \textit{Structural missingness:} drug use is only available for ages 20--59,
        and chronic conditions, smoking and marital status only from age 20 (step 1).
        We add missing-indicators rather than imputing these as ``no''.
\end{itemize}

\textbf{Step 3 --- Exploratory analysis.} Report prevalence by cycle, sex, age and
ethnicity with the NHANES survey weights (\var{WTMEC2YR} combined over cycles, strata
\var{SDMVSTRA}, PSU \var{SDMVPSU}) so the figures describe the US adult population.
Produce univariate risk ratios and a correlation check for collinear features (e.g.\
antidepressant use vs.\ mental-health visit).

\textbf{Step 4 --- Modelling.} Temporal split: \textbf{train 2005--2014 (26,057
adults, 983 cases) / test 2015--2018 (10,202 adults, 383 cases)}. Hyperparameters are
tuned with stratified 5-fold cross-validation inside the training set only. We fit
four nested feature sets (demographics $\rightarrow$ routine data $\rightarrow$ PHQ-8
only $\rightarrow$ routine + PHQ-8) with two algorithms:
\begin{itemize}
  \item \textbf{L2- or elastic-net logistic regression}, interpretable as odds ratios;
  \item \textbf{gradient boosting} (scikit-learn \var{HistGradientBoostingClassifier}
        or XGBoost).
\end{itemize}
The imbalance comparison (sub-question 1, H4) runs each model three ways: with no
correction, with balanced class weights, and with SMOTE applied inside the CV folds only.

\textbf{Step 5 --- Evaluation.}
\begin{itemize}
  \item Primary metric: \textbf{PR-AUC}, compared against the prevalence baseline.
        Secondary: ROC-AUC, Brier score and a calibration plot.
  \item Clinical operating point: \textbf{recall and precision when flagging the top
        5\% of patients}, reported as ``flag 1 in 20 patients $\rightarrow$ catch X\%
        of cases, of whom Y\% are true positives''.
  \item Hypothesis tests via bootstrap CIs on the test set.
  \item Fairness (H5): ROC-AUC, PR-AUC relative to subgroup prevalence, and recall at
        top 5\% by sex, age group and ethnicity.
\end{itemize}

\textbf{Step 6 --- Explainability.} SHAP values for the gradient-boosting model and
odds ratios for logistic regression. Global importance ranking plus 2--3 individual
patient explanations (e.g.\ ``flagged mainly because of antipsychotic prescription,
poor self-rated health and not working'').

\textbf{Tools:} Python (pandas, scikit-learn, imbalanced-learn, shap, matplotlib); code
in a shared Git repository.

% =================================================================== 5
\section{Ethics and limitations}

\begin{itemize}
  \item \textbf{Public, de-identified data.} NHANES is released for public research
        and needs no ethics approval. We report only aggregate results.
  \item \textbf{Not a diagnostic tool.} The outcome is \textit{self-reported
        ideation}, not suicide attempts or deaths, and the design is
        \textbf{cross-sectional}. The model \textit{flags}; it does not \textit{predict
        the future}. We discuss the harm of false positives (stigma, unnecessary
        referrals) and false negatives.
  \item \textbf{Transferability.} The data are from the US; Swiss prescribing,
        insurance and care pathways differ. BFS figures are used for context only.
  \item \textbf{Label quality.} PHQ item 9 captures passive thoughts as well as active
        ideation, so it is known to over-identify risk~\cite{kroenke2001}.
\end{itemize}

% =================================================================== 6
\section{Out of scope (future work)}

NLP on clinical notes (ClinicalBERT), multimodal late fusion, and 30-day
post-discharge prediction. These need credentialed access to MIMIC-IV/ScAN and
longitudinal outcome data, and they are the natural next step once the
structured-data signal has been quantified here.

% =================================================================== 7
\section{Timeline and task split}

\begin{table}[h]
\centering
\small
\begin{tabularx}{\textwidth}{@{}l L l@{}}
\toprule
\textbf{Period} & \textbf{Milestone} & \textbf{Owner} \\
\midrule
25.09 -- 01.10 & Download, merge and harmonise all cycles; outcome definition;
                 missingness report & 2 members \\
25.09 -- 09.10 & Write preliminary study (this plan, max.\ 2 pages + appendix) & all \\
\textbf{09.10} & \textbf{Hand-in preliminary study} & \\
09.10 -- 16.10 & Feature engineering (medication mapping, multimorbidity), weighted EDA
               & 2 members \\
12.10 -- 21.10 & Modelling, imbalance comparison, bootstrap tests & 2 members \\
19.10 -- 24.10 & SHAP, fairness, calibration, clinical operating point & 1--2 members \\
22.10 -- 28.10 & Slides ($\le$ 10), code clean-up, GenAI disclosure & all \\
\textbf{28.10} & \textbf{Hand-in final project (.zip)} $\cdot$ 30.10 presentation in
                 Lucerne & \\
\bottomrule
\end{tabularx}
\end{table}

% =================================================================== references
\bibliographystyle{unsrtnat}
\bibliography{references}

% =================================================================== appendix
\appendix
\section{Harmonisation rules (step 1)}
\label{app:harmonisation}

\begin{itemize}
  \item \textit{``Refused / don't know'' codes depend on the variable} (7/9, 77/99 or
        none) and are recoded per variable, because 7 and 9 are real answers for some
        variables, such as hours of sleep.
  \item \textit{Healthcare visits:} \var{HUQ050} (2005--2012) $\rightarrow$
        \var{HUQ051} (2013+). The finer 2013 categories nest exactly into the earlier
        ones and are collapsed to them. The 2013 question no longer counts
        emergency-room and home visits, a definitional break we report.
  \item \textit{Sleep hours:} \var{SLD010H} (whole hours, 12 = ``12 or more'')
        $\rightarrow$ \var{SLD012} in the new 2015 sleep module (half hours, up to
        14.5); both are clipped to 2--12\,h.
  \item \textit{Age} is top-coded at 85 in 2005--06 and at 80 afterwards; we use 80
        throughout.
  \item \textit{Education} for 18--19-year-olds comes from the youth question
        \var{DMDEDUC3}, because the adult question starts at 20.
  \item \textit{High cholesterol} was asked until 2010 only if cholesterol had been
        checked; people never checked are coded ``no'', as everyone is asked from 2011.
  \item \textit{Age coverage} differs between cycles: smoking and cholesterol are asked
        from 20 until 2012 and from 18 afterwards, hard-drug use covers ages 20--59 in
        2005--06, 20--69 in 2007--08 and 18--69 from 2009, and marital status covers
        14+ in 2005--06 and 20+ later. Each variable is kept only for the ages asked
        in every cycle (20+, drug use 20--59), so missingness does not reveal the
        survey year in the temporal validation (H5).
\end{itemize}

The full variable list with NHANES sources per cycle is generated with the code
(\var{docs/codebook.md} in the project repository).

\end{document}
